% !TeX root = ../papers/nc-p3.tex
\section{平面曲線の一回のブローアップと有限な反対階数八の全次数判定}
\label{sec:finite-rank-eight-all-marks}
\begin{lemma}\label{lem:finite-carrier-blowup}
有限窓の共通印付き四空間の階数一切断が長さ3とする。
印の平面局所環を $P=C/(g)$、$C$ を重複度 $d$ の平面超曲面完備局所環、
$g$ を極大イデアルの二乗に入る非零因子とする。
$d=2$ なら対称性判定が成立する。切断が非Gorensteinなら $d=3$ でも成立する。
曲線の被約性は仮定しない。
\end{lemma}
\begin{proof}
$C$ の一般の一次パラメータを $x$、もう一つの座標を $y$ とする。
Weierstrass表示を $y^d+\sum_{j=1}^d a_j(x)y^{d-j}=0$ とすれば
$a_j\in(x^j)$ である（準備で除いた因子は単元なので位数を保つ）。
$C'=C[y/x]\subset C[1/x]$ は $1,z,\ldots,z^{d-1}$, $z=y/x$ を
$k[[x]]$ 基底に持ち、$x^{d-1}C'\subset C$、$\mathfrak m_C C'=x C'$。
実際 $x^d$ で割ったモニック方程式を使え、線型独立性は $x$ の逆元を加えて従う。
$h=g/x\in xC'$ で $C'/(h)$ はArtin環である。

共通行列を $A$、規格化したもう一つの行列との差を $Z$ とし、行・列核の
三座標を $l,r$ とする。全てを $C$ に持ち上げ、差は固定三空間 $N$ 内で取る。
$l,r$ の成分はどちらも $\mathfrak m_C$ を生成する。実際 $P$ で極大イデアルを生成し、
$g\in\mathfrak m_C^2$ なので一次部が基底を含む。
従って $L=l/x$, $R=r/x$ は $C'$ 上で各々単位イデアルを生成する。
$A,Z$ は $h$ を法として $L,R$ を共通核に持つ。これは元の縮約が $g$ の倍であるためである。

$C'/(h)$ の各局所因子で、剰余体でのこの二核を消す $E$ の空間は次元2以下である。
実際三次元なら二行二列行列の射影平面に階数1の曲線ができる。
$A$ は固定超平面 $N$ の外の単元座標を持つので、$N$ 上の縮約核は剰余体で高々一次元。
二つの主元で消去すれば、差は零、または $Z=eV$ で $V$ は剰余体で $A$ と独立である。
後者で $\operatorname{ann}e$ を法とすれば $A,V$ は共通核を持つ相対割線を与える。
この商でも $L,R$ は射影座標となり、剰余体で階数1集合は有限なので
行列式が割線上に切る相対次数2の部分は平坦である。

Gorenstein三重点では定理\ref{thm:finite-marked-artin-secants}の三方向上で
消える任意の二次式 $Q$ が $Q(L)Z\in h C'$ を成分ごとに満たす。
従って $Q(l)Z\in xg C'$ であり、$d=2$ なら $Q(l)Z=0$ が $P$ 上で成り立つ。
印の外での同じ消滅と三方向のイデアル列を用いる同定理の証明で対称性が従う。

非Gorenstein三重点では定理\ref{thm:finite-marked-planar-triple}の $E_D$ 座標を用いる。
$L=(L_0,L')$ とすると、相対割線の左核には $L_0=0$ が必要なので $L_0 Z\in h C'$。
$A=\bigl(\begin{smallmatrix}a&b^t\\b&tD\end{smallmatrix}\bigr)$ では
$a$ が単元、$b$ と $Z$ の全成分は印で消え $x C'$ に入る。
各局所因子で $D^t L'$ の少なくとも一成分は単元である。
実際、$A$ の剰余は $E_{11}$ の非零倍であり、$L$ の剰余は非零でこれを消す。
左核の下二式 $tD^t L'+b L_0\in h C'$ に $Z$ の各成分を掛けると
$b L_0 Z\in xh C'$ かつ $hZ\in xh C'$ である。可逆な成分で消去して
\[
                 t Z\in g C',\qquad l_0 Z=x L_0 Z\in g C'
\]
を得る。ここで単元での消去は $C'$ の各極大イデアルで局所化して行える。
$d\le3$ なら $x^2 C'\subset C$ により $\mathfrak m_P^2 t Z=\mathfrak m_P^2 l_0 Z=0$。
第二辺の評価 $f$ は局所的に $Z$ の成分の線型結合であり、印の外では既出の相対割線の
議論により $tf,l_0f$ は零である。印で極大イデアルの二乗が消す各元の支持環は
Gorensteinで長さ2以下、対応する隣の一次座標が全射する。
隣接座標積の完全跡直交から $tf=l_0 f=0$ が全域で成立する。
以後は長さ6の二平面切断の証明を適用できる。
すなわち全評価の空間は高々一次元で、非零ならその一生成元の支持は二平面切断全体である。
定理\ref{thm:finite-nodal-cuspidal-marks}の証明の通り、平面印におけるこの生成元の
支持環は曲線状であり（この段に曲線 $C$ の条件は要らない）、
長さ5以上で消える余因子と長さ6の二次補間が矛盾する。
\end{proof}

\begin{theorem}\label{thm:finite-rank-eight-all-marks}
四辺が有限な反対階数8の標準有限Frobenius窓は、全て次の有限判定に入る。
中央支持3の点があれば全次数で正則である。それ以外では少なくとも一方の反対対称性が成立する。
その対称核が階数2以下の非零元を含めば全次数で正則であり、残りでは
両対称性と定理\ref{thm:bisymmetric-two-length-test}の条件が正則性に必要十分である。
\end{theorem}
\begin{proof}
中央支持3の認証は定理\ref{thm:finite-cubic-separators}による。
曲線状の印・伸長型の印・短い局所環の場合を除けば、残る印は全て平面局所完全交叉で
閉点数は高々2である。両側の共通四空間は階数8で自動的に存在する。
未判定なら定理\ref{thm:high-order-octad-bisymmetry}により、印における
二定義式の最小位数 $d$ は2または3である。位数 $d$ の式で曲線 $C$ を取り、
他方をその上の非零因子 $g$ とする。

一つの印を取る。いずれかの側の階数一切断が正次元または長さ2以下なら既出の対称性を使う。
長さ4では一点支持なら系\ref{cor:single-support-quartic-slice}を使い、
二点支持で核側が非対称なら系\ref{cor:two-support-quartic-discrepancy}が同じ印に
重平面因子を持つ別の分離式を与える。その切断は定理\ref{thm:finite-split-separator-slices}より
非Gorensteinの長さ3である。従って残る場合にはこの印で長さ3の切断を使える。
$d=2$ なら前補題を適用する。$d=3$ では局所定義式に二次の初期式がない。
定理\ref{thm:finite-high-order-common-octad}の最初のブロック計算
$C_1(x)Rx=(Lx)^t C_1(x)=0$ はまさにこの次数までしか使わず、
階数一切断の印での接次元が2以上となる。長さ3なら平方零の非Gorenstein局所環なので、
再び前補題が適用できる。これで少なくとも片側が対称になり、
系\ref{cor:all-defective-one-sym-finite}が述べた判定を与える。
\end{proof}
